\documentclass[conference]{IEEEtran}
\IEEEoverridecommandlockouts
% ---------------------------------------------------------------------------
% Este arquivo segue o "IEEE Conference Template" oficial do Overleaf, como
% exigido pela especificação da disciplina. Para compilar: crie um projeto
% novo no Overleaf a partir desse template, substitua o main.tex pelo
% conteúdo deste arquivo e faça upload de bench/fps_vs_resolucao.png na
% pasta do projeto (referenciada na Figura 1, Seção IV).
% ---------------------------------------------------------------------------
\usepackage[utf8]{inputenc}
\usepackage[T1]{fontenc}
\usepackage[brazil]{babel}
\usepackage{cite}
\usepackage{amsmath}
\usepackage{amssymb}
\usepackage{graphicx}
\usepackage{textcomp}
\usepackage{xcolor}
\usepackage{booktabs}
\usepackage{url}

\def\BibTeX{{\rm B\kern-.05em{\sc i\kern-.025em b}\kern-.08em
    T\kern-.1667em\lower.7ex\hbox{E}\kern-.125emX}}

\begin{document}

\title{Visualização Interativa de Elevação Costeira com Nível de Detalhe
Adaptativo e Simulação Oceânica para Análise de Cenários de Subida do
Nível do Mar}

\author{
\IEEEauthorblockN{Guilherme Augusto Scaglia, João Pedro Denardo, Pedro Henrique Oliveira de Souza}
\IEEEauthorblockA{
Engenharia de Computação\\
Fundação Hermínio Ometto (FHO)\\
Araras, SP, Brasil\\
\{scaglia, denardo749, pedro1204\}@alunos.fho.edu.br}
}

\maketitle

\begin{abstract}
Este trabalho apresenta um sistema interativo de visualização científica
para a análise de cenários de erosão costeira e subida do nível do mar,
implementado sobre um Modelo Digital de Elevação (DEM) real da Baixada
Santista (SP), obtido por meio do serviço \emph{AWS Terrain Tiles}, cujos
dados primários combinam observações de radar orbital (SRTM) e batimetria,
isto é, a profundidade do fundo do mar (GEBCO). A hipótese central deste trabalho é a de que o nível de detalhe
adaptativo (LOD) não constitui, isoladamente, o objetivo da pesquisa, mas
sim o mecanismo viabilizador de uma ferramenta de visualização de
inundação costeira em tempo real. Buscou-se determinar experimentalmente
a partir de qual carga geométrica a técnica CDLOD (\emph{Continuous
Distance-Dependent Level of Detail}), de Strugar, torna-se indispensável
para a manutenção da interatividade da aplicação, tendo dispositivos
móveis como público-alvo final. A superfície oceânica é modelada por seis
componentes de ondas de Gerstner, na formulação de Fournier e Reeves, com
altura de base ajustável em tempo real, o que permite a visualização de
cenários de subida do mar de 0, +0,5, +1 e +2 metros sobre o relevo
costeiro. Um teste de desempenho (\emph{benchmark}) automatizado, executado sem
janela gráfica em uma GPU NVIDIA RTX 3060, comparou a renderização com e sem a técnica de LOD
sobre cinco resoluções de terreno, identificando um ponto de cruzamento
(ponto de virada) em aproximadamente cinco milhões de triângulos, a
partir do qual o custo do CDLOD passa a ser inferior ao da malha completa.
O sistema foi desenvolvido em linguagem Python, com as bibliotecas
ModernGL, GLFW e PyGLM.
\end{abstract}

\begin{IEEEkeywords}
computação gráfica, visualização científica, nível de detalhe,
renderização de terreno, simulação de oceano, erosão costeira, subida do
nível do mar
\end{IEEEkeywords}

\section{Introdução}

O aumento do nível médio dos oceanos, impulsionado por mudanças
climáticas de origem antrópica, constitui atualmente um dos fenômenos
ambientais de maior impacto sobre regiões costeiras em escala global.
Esse processo intensifica a erosão de linhas de costa, amplia a
vulnerabilidade de áreas urbanas litorâneas a eventos de inundação e
representa um risco direto a infraestruturas, ecossistemas costeiros e
populações que habitam essas regiões. A previsão e a comunicação visual
desses riscos constituem, portanto, elementos centrais tanto para o
planejamento urbano e ambiental quanto para a conscientização pública
sobre o tema.

Conforme discutido por Brown \emph{et al.} \cite{brown2006}, um dos
principais entraves na comunicação de risco costeiro à população e a
gestores públicos é a baixa capacidade das ferramentas tradicionais de
geovisualização em transmitir a dinâmica real de um fenômeno costeiro,
uma vez que a maioria delas adota o denominado modelo de ``banheira''
(\emph{bathtub model}), no qual toda área abaixo de uma determinada cota
de elevação é simplesmente marcada como inundada, sem qualquer simulação
da superfície da água. Para mitigar essa limitação, o presente trabalho
fundamenta-se na integração entre a renderização de terreno em tempo real
e a simulação procedural de ondas oceânicas sobre dados reais de relevo
(topografia) e de profundidade do fundo do mar (batimetria).

Cabe demarcar, contudo, o posicionamento exato da contribuição técnica
deste trabalho. Diferentemente de um projeto cujo objetivo fosse propor
uma nova técnica de nível de detalhe, este trabalho trata o LOD adaptativo
como um meio, não como um fim. A pergunta de pesquisa que orienta a
avaliação experimental (Seção IV) não é ``o CDLOD funciona?'' --- resposta
já estabelecida por Strugar \cite{strugar2009} --- mas sim ``a partir de
que carga geométrica, num cenário real de visualização de inundação
costeira, o CDLOD deixa de ser um refinamento de desempenho e passa a ser
uma condição necessária para a interatividade da aplicação?''. Essa
pergunta ganha relevância particular quando o público-alvo final da
ferramenta inclui dispositivos móveis, cuja capacidade de processamento
gráfico é substancialmente inferior à de uma GPU de \emph{desktop}.

Diferentemente de sistemas clássicos de geovisualização baseados em
Sistemas de Informação Geográfica (SIG), que representam a subida do mar
por meio de camadas estáticas sobrepostas a um mapa bidimensional, este
projeto propõe uma ruptura na forma de apresentação do fenômeno para o
público leigo e para gestores ambientais: a construção de uma cena
tridimensional interativa, sobre um DEM real, na qual o observador pode
navegar livremente pelo relevo costeiro e ajustar, em tempo real, a
altura da superfície oceânica simulada, observando diretamente quais
áreas seriam submersas em cada cenário de subida do nível do mar.

O sistema foi desenvolvido em Python, com a biblioteca ModernGL provendo
acesso à API gráfica OpenGL, a biblioteca GLFW gerenciando janela e
entrada, e a biblioteca PyGLM fornecendo as operações de álgebra linear
associadas às transformações de câmera e projeção. A solução não pretende
apenas renderizar um terreno esteticamente convincente, mas também
constituir uma ferramenta de análise visual de baixo custo computacional,
complementar às abordagens tradicionais baseadas em SIG.

\subsection{Justificativa da Escolha do Cenário e das Técnicas}

A escolha do cenário de erosão costeira, em detrimento de outras
aplicações possíveis para as técnicas de renderização de terreno,
justifica-se por três razões. Em primeiro lugar, trata-se de um cenário
que já possui embasamento em trabalhos de geovisualização científica
\cite{brown2006,dipaola2023}. Em segundo lugar, o parâmetro central da
simulação de água (a altura de base ajustável) corresponde diretamente a
uma variável de interesse real: as projeções de subida do mar, expressas
em metros. Em terceiro lugar, o problema impõe uma restrição de escala
compatível com o escopo do projeto: diferentemente da renderização de um
planeta inteiro \cite{clasenhege2006}, um recorte costeiro de dimensão
local --- a Baixada Santista, detalhada na Seção III-A --- permite
simplificar a topologia do terreno de esférica para plana, sem
descaracterizar o problema.

A escolha específica das técnicas --- CDLOD para o terreno e ondas de
Gerstner para a água --- segue o mesmo raciocínio de custo-benefício,
posteriormente confirmado de forma experimental na Seção IV: o CDLOD
dispensa estruturas de dados hierárquicas mais complexas do que outras
técnicas de LOD, e o modelo de Gerstner dispensa o custo computacional
das transformadas de Fourier exigidas por modelos espectrais de oceano
(como o de Tessendorf), mantendo fidelidade visual suficiente para o
objetivo de comunicação de risco proposto, sem inviabilizar a execução em
hardware de menor capacidade de processamento --- objetivo final da
aplicação, tendo em vista o público-alvo de dispositivos móveis.

\section{Fundamentação Teórica e Trabalhos Relacionados}

A revisão bibliográfica abrange sete trabalhos, organizados em três
frentes: a renderização eficiente de terrenos extensos, a simulação de
ondas oceânicas e a geovisualização de risco costeiro. As subseções
II-A a II-D descrevem o que cada frente propõe e quais são suas limitações
para o problema tratado; a Subseção II-E reúne essas observações e
explicita o que este trabalho adota de cada uma e no que se diferencia.

\subsection{Renderização de Terreno em Larga Escala}

Um terreno real extenso, representado como malha de triângulos, pode
conter milhões de triângulos, o que torna inviável desenhá-lo por inteiro
a cada quadro. A solução usual é o \emph{nível de detalhe} (LOD, do
inglês \emph{level of detail}): usar muitos triângulos nas regiões
próximas ao observador, onde o detalhe é perceptível, e poucos nas
regiões distantes.

Tanner \emph{et al.} \cite{tanner1998} introduziram o \emph{clipmap}, uma
estrutura que mantém em memória uma ``janela'' de dados em alta
resolução centrada no observador, cercada por janelas sucessivamente
maiores e de resolução menor. Assim, apenas os dados próximos precisam
ser armazenados em alta resolução, o que permite renderizar mapas de
elevação maiores do que a memória da GPU comportaria. Clasen e Hege
\cite{clasenhege2006} estenderam essa ideia a terrenos esféricos, com
anéis circulares de detalhe, de modo a visualizar planetas inteiros a
partir de modelos digitais de elevação de alta resolução.

Ambos os trabalhos visam grandes extensões, até a escala planetária. Para
um recorte costeiro local, de superfície aproximadamente plana, a
parametrização esférica e a manutenção de janelas aninhadas de dados
acrescentam complexidade que o problema não exige.

\subsection{Nível de Detalhe Contínuo Baseado em Distância}

Strugar \cite{strugar2009} propôs o CDLOD (\emph{Continuous
Distance-Dependent Level of Detail}). O terreno é organizado em uma
\emph{quadtree}, estrutura em árvore na qual um bloco quadrado pode ser
dividido em quatro blocos menores, e assim recursivamente. Blocos
próximos ao observador são subdivididos, ganhando detalhe; blocos
distantes permanecem grandes e grosseiros. A decisão depende de uma
função contínua da distância, e a troca de nível é suavizada por
\emph{geomorphing}: os vértices deslizam gradualmente de uma resolução
para a seguinte, evitando saltos visuais abruptos e fendas entre blocos
vizinhos.

Em comparação com o \emph{clipmap}, o CDLOD é mais simples de
implementar, pois dispensa janelas de dados aninhadas e se resume a
escolher, a cada quadro, quais blocos desenhar. O trabalho original foi
proposto e avaliado para terrenos genéricos, sem vínculo com um cenário
de aplicação específico.

\subsection{Simulação Procedural de Ondas Oceânicas}

Fournier e Reeves \cite{fournierreeves1986} propuseram um dos primeiros
modelos procedurais de ondas para síntese de imagens, baseado no modelo
de Gerstner. Nele, cada ponto da superfície da água descreve uma
trajetória circular ou elíptica, o que produz cristas estreitas e vales
largos, como em ondas reais; somando várias ondas de tamanhos
diferentes, obtém-se uma superfície de aspecto natural, a um custo muito
inferior ao dos modelos espectrais. O trabalho considera ainda a
influência da batimetria (profundidade do fundo do mar) sobre a
refração e a arrebentação das ondas próximas à costa. Welch \cite{welch}
descreve a implementação de uma simulação de oceano em tempo real, com
ênfase no desempenho interativo em GPU.

Esses trabalhos modelam o oceano com nível médio fixo; o nível da água
não é tratado como variável de cenário.

\subsection{Geovisualização de Erosão Costeira e Subida do Nível do Mar}

Brown \emph{et al.} \cite{brown2006} discutem a importância de visualizar
linhas de costa em evolução como ferramenta de conscientização e
planejamento frente à erosão, e propõem uma metodologia que integra a
modelagem científica da erosão a modelos digitais de elevação (DEM) e a
sistemas de informação geográfica. De forma semelhante, DiPaola,
Bhardwaj e Sam \cite{dipaola2023} constroem mapas interativos de subida
do nível do mar sobre dados topográficos reais (LiDAR aéreo de alta
resolução) do North Shore de Vancouver, Canadá, quantificando a
porcentagem de área afetada em diferentes cenários.

Esses trabalhos são cientificamente rigorosos, mas apresentam o
fenômeno por meio de mapas bidimensionais ou modelos de inundação
estáticos, sem simular a superfície da água.

\subsection{Síntese e Posicionamento deste Trabalho}

Da primeira frente, retivemos a ideia de concentrar o detalhe perto do
observador. Como o recorte de interesse é local e plano, adotamos a
\emph{quadtree} do CDLOD em lugar da parametrização esférica do
\emph{clipmap}. Da segunda frente, adotamos o modelo de Gerstner, na
formulação de Fournier e Reeves, com uma diferença: a altura média da
água é um parâmetro ajustável em tempo real, usado para representar
diferentes cenários de subida do mar. A terceira frente fornece a
motivação e a escolha do cenário; diferimos dela ao oferecer uma cena
tridimensional navegável, com água animada, em vez de mapas
bidimensionais estáticos, sob a hipótese de que a dinâmica visual da
água reforça a percepção de risco por parte do observador.

\begin{table}[t]
\caption{Comparação qualitativa com trabalhos relacionados}
\label{tab:comparacao}
\centering
\begin{tabular}{@{}p{2.0cm}p{2.2cm}p{1.7cm}p{1.6cm}@{}}
\toprule
\textbf{Trabalho} & \textbf{Técnica de LOD} & \textbf{Água dinâmica} & \textbf{Contexto} \\
\midrule
Clasen e Hege \cite{clasenhege2006} & Clipmap esférico & Não & Planeta \\
Strugar \cite{strugar2009} & CDLOD (quadtree) & Não & Terreno \\
DiPaola \cite{dipaola2023} & N/A (SIG) & Não & Costa real \\
\textbf{Proposta} & CDLOD adaptado & Sim & Costa (erosão) \\
\bottomrule
\end{tabular}
\end{table}

Conforme a Tabela \ref{tab:comparacao}, a contribuição deste trabalho
não é uma técnica inédita, mas a integração de técnicas de renderização
de terreno e de simulação de água --- em geral aplicadas a jogos ou a
visualização genérica --- a um cenário de análise ambiental, aliada à
determinação experimental do ponto a partir do qual o LOD se torna
necessário (Seção IV).

\section{Metodologia e Modelagem Computacional}

A metodologia deste trabalho foi estruturada em um \emph{pipeline}
modular de cinco estágios: (A) aquisição e geração dos dados de elevação;
(B) estruturação do terreno em quadtree e nível de detalhe; (C) modelo de
câmera de navegação; (D) simulação da superfície oceânica; e (E)
ambiente de desenvolvimento. A avaliação experimental de desempenho, cujo
protocolo decorre diretamente desta metodologia, é apresentada em seção
própria (Seção IV).

\subsection{Aquisição e Geração dos Dados de Elevação}

A área de estudo compreende a Baixada Santista (municípios de Santos, São
Vicente e Guarujá, SP), delimitada pela caixa geográfica de latitude
$-23{,}70$ a $-24{,}10$ e longitude $-46{,}60$ a $-46{,}05$, selecionada
por reunir, em um único recorte, relevo montanhoso (Serra do Mar),
planície costeira baixa e o sistema estuarino de Santos e São Vicente (região onde rios desembocam
no mar).

O relevo é representado por um mapa de elevação (\emph{heightmap}): uma
matriz $H \in \mathbb{R}^{n \times n}$, com $n = 1024$, na qual cada
célula $H(x,z)$ armazena uma altura normalizada. Consideram-se dois
métodos de obtenção de $H$, sendo o primeiro a fonte primária da
implementação de referência.

\textbf{Dados reais (fonte primária).} O terreno é obtido a partir de
\emph{tiles} de elevação do serviço \emph{AWS Terrain Tiles}, distribuído
publicamente sem exigência de chave de API, no formato \emph{terrarium},
que codifica a altitude real em metros diretamente nos três canais de cor
de uma imagem PNG, segundo:
\begin{equation}
h_{\text{real}} = (R \cdot 256 + G + B/256) - 32768
\label{eq:terrarium}
\end{equation}
em que $R$, $G$ e $B$ são os valores dos canais de cor do \emph{tile}
(0--255) e $h_{\text{real}}$ é a elevação decodificada, em metros. Os
dados subjacentes combinam observações orbitais de radar (SRTM, para a
porção emersa) e batimetria GEBCO (para o leito oceânico), permitindo
representar simultaneamente relevo terrestre e profundidade marinha em
uma única matriz contínua.

Em razão de os \emph{tiles} brutos apresentarem ruído pontual
característico de dados de sensoriamento remoto, aplica-se um
pós-processamento composto por: (i) filtro de mediana $3 \times 3$, em
duas passadas, para a remoção de \emph{outliers} isolados
(\emph{despiking}) sem suavização excessiva das feições do relevo; (ii)
reamostragem por interpolação bilinear para a resolução-alvo de
$1024 \times 1024$ células; e (iii) imposição de um piso mínimo de
elevação em $-40$ m, o que evita que artefatos de batimetria produzam
profundidades irrealistas na região de interesse.

A elevação real, em metros, é convertida para unidades de mundo (u.m.) do
motor gráfico por meio de um fator de escala horizontal e de um fator de
exagero vertical, aplicados conjuntamente:
\begin{equation}
y_{\text{mundo}} = \frac{h_{\text{real}}}{\sigma} \cdot E
\label{eq:escala}
\end{equation}
em que $\sigma = 167$ m/u.m. é a escala do mundo (um terreno de 300 u.m.
de lado corresponde a aproximadamente 50 km reais, cobrindo a extensão da
área de estudo) e $E = 6$ é o fator de exagero vertical. O exagero
vertical é uma prática consolidada em geovisualização, necessária neste
caso porque o relevo real da Serra do Mar (amplitude aproximada de
1,2 km) representa uma fração muito pequena da extensão horizontal da
área de estudo (aproximadamente 50 km): sem o fator $E$, o relevo
apareceria visualmente quase plano.

\textbf{Geração procedural (modo alternativo).} Como alternativa
independente de conectividade de rede, contempla-se a geração de $H$ por
ruído coerente do tipo \emph{Simplex noise}, somado em múltiplas oitavas
(\emph{fractal Brownian motion}):
\begin{equation}
H(x,z) = \frac{1}{\sum_{i=0}^{k-1} p^i} \sum_{i=0}^{k-1} p^i \cdot
\text{Simplex}(f_0 2^i x,\, f_0 2^i z)
\label{eq:fbm}
\end{equation}
com $k = 5$ oitavas e persistência $p = 0{,}5$.

Independentemente da origem dos dados, cada célula de $H$ é utilizada
para o deslocamento vertical dos vértices de uma grade regular
(\emph{displacement mapping}), transferida à GPU como uma textura de
ponto flutuante de canal único (R32F) --- um valor de altura por texel,
lido diretamente pelo \emph{vertex shader} no estágio descrito na
Subseção III-B.

\subsection{Estruturação do Terreno em Quadtree e Nível de Detalhe (CDLOD)}

O algoritmo CDLOD \cite{strugar2009} resolve o seguinte problema:
renderizar toda a extensão do terreno na resolução máxima é
computacionalmente inviável em tempo real, mas reduzir a resolução de
maneira uniforme degrada a qualidade visual nas regiões próximas ao
observador, onde o detalhe é mais perceptível. A solução consiste em
variar a densidade de triângulos em função da distância à câmera.

\textbf{Estrutura de dados.} O terreno é particionado recursivamente em
uma quadtree: o nó raiz cobre toda a extensão do terreno; cada nó filho
cobre um quadrante da área do nó pai, até uma profundidade máxima
$L = 5$ (valor implementado, suficiente para representar a região
costeira de interesse com resolução submétrica nas proximidades do
observador). Cada nó armazena a mesma quantidade fixa de vértices,
instanciada como uma grade regular de $N \times N$ pontos ($N = 33$,
correspondendo a $32 \times 32$ subdivisões, ou 2048 triângulos por nó),
de modo que um nó em profundidade $\ell$ cobre uma área
$\left(\frac{W}{2^{\ell}}\right)^2$, sendo $W$ a largura total do
terreno. Nós mais profundos representam, portanto, áreas menores com a
mesma densidade de vértices, resultando em maior resolução por unidade de
área.

\textbf{Critério de seleção de nível.} Para cada nó de lado $\tau$
(tamanho do bloco na profundidade correspondente), calcula-se a distância
euclidiana $d$ entre a posição da câmera $\vec{c}$ e o ponto mais próximo
da borda da caixa delimitadora do nó. O critério de subdivisão
implementado é:
\begin{equation}
\text{subdividir se } d < \tau \cdot 2{,}2
\label{eq:subdivisao}
\end{equation}
ou seja, um nó é subdividido em quatro filhos de maior detalhe sempre que
a câmera estiver mais próxima de sua borda do que 2,2 vezes o próprio
tamanho do bloco; caso contrário, o nó é renderizado em sua resolução
atual. Como o tamanho do bloco $\tau$ é reduzido à metade a cada nível de
profundidade, esse critério produz um efeito equivalente a um raio de
comutação que dobra a cada nível --- analogamente a um esquema de anéis
concêntricos de detalhe centrados na câmera --- sem exigir o
armazenamento explícito de distâncias de corte por nível. Na
implementação de referência, esse critério resulta tipicamente em 200 a
230 blocos ativos simultaneamente, frente aos $4^5 = 1024$ blocos-folha
que existiriam caso a árvore inteira fosse forçada à profundidade máxima.

\textbf{Eliminação de descontinuidades (\emph{cracks}).} Blocos
adjacentes renderizados em níveis de detalhe diferentes possuem
densidades de vértices distintas em sua borda compartilhada, o que pode
gerar fendas visuais (\emph{cracks}) caso os vértices das duas malhas não
coincidam exatamente na fronteira. O CDLOD trata esse problema por meio
de uma técnica de \emph{geomorphing}: na faixa de transição de cada
bloco, os vértices de índice ímpar da grade deslizam progressivamente em
direção à posição dos vértices de índice par correspondentes, eliminando
tanto as fendas quanto o efeito de salto abrupto (\emph{popping}) quando
a câmera se move e um bloco muda de nível de detalhe.

\subsection{Modelo de Câmera de Navegação Livre}

A navegação pela cena é realizada por uma câmera de movimento livre
(\emph{free-fly camera}), cuja orientação é parametrizada pelos ângulos
de Euler \emph{yaw} ($\psi$, rotação em torno do eixo vertical) e
\emph{pitch} ($\theta$, rotação em torno do eixo horizontal lateral). O
vetor de direção para onde a câmera aponta (\emph{front}) é derivado
desses ângulos por:
\begin{equation}
\vec{f} = \big(\cos\psi\cos\theta,\ \sin\theta,\ \sin\psi\cos\theta\big)
\label{eq:front}
\end{equation}

A partir de $\vec{f}$ e do vetor global ``para cima''
$\vec{u}_{\text{mundo}} = (0,1,0)$, calculam-se os vetores ortogonais
\emph{right} ($\vec{r} = \vec{f} \times \vec{u}_{\text{mundo}}$,
normalizado) e \emph{up} ($\vec{u} = \vec{r} \times \vec{f}$), que formam
a base local da câmera. Esses três vetores, associados à posição
$\vec{c}$, constroem a matriz de visualização (\emph{view matrix}) por
meio da transformação \emph{look-at}, que expressa a cena no referencial
da câmera. A cada quadro, entradas de teclado incrementam $\vec{c}$ na
direção de $\vec{f}$, $\vec{r}$ ou $\vec{u}$, segundo uma velocidade de
translação de 40 unidades de mundo por segundo em modo normal e 120
unidades de mundo por segundo em modo acelerado (\emph{turbo}), ambas
escaladas pelo tempo decorrido entre quadros (\emph{delta time}) para
garantir independência da taxa de quadros; o deslocamento do mouse
incrementa $\psi$ e $\theta$, com $\theta$ restrito ao intervalo
$(-89^{\circ}, 89^{\circ})$ para evitar a inversão do eixo vertical da
câmera (\emph{gimbal flip}). A matriz de projeção perspectiva é
parametrizada com campo de visão vertical de $60^{\circ}$ e planos de
recorte próximo e distante em 0,1 e 2000 unidades de mundo,
respectivamente --- intervalo dimensionado para conter integralmente a
extensão do terreno descrita na Subseção III-A sem introduzir imprecisão
numérica no \emph{depth buffer}.

Além de sua função de navegação, $\vec{c}$ constitui a entrada direta do
critério de distância descrito na Subseção III-B: a cada quadro, a mesma
posição de câmera utilizada para construir a matriz de visualização é
reutilizada para o cálculo de $d$ em relação a cada nó da quadtree do
terreno, o que evita o cálculo redundante de duas representações
distintas da posição do observador.

\subsection{Simulação da Superfície Oceânica}

A superfície da água é modelada pelo modelo de ondas de Gerstner (também
chamado trocoidal), na formulação de Fournier e Reeves
\cite{fournierreeves1986}. Diferentemente de uma onda senoidal simples,
que desloca os pontos da superfície apenas verticalmente, o modelo de
Gerstner desloca cada ponto tanto na vertical quanto na direção
horizontal de propagação da onda, de modo que partículas da superfície
descrevem trajetórias elípticas --- o que produz o perfil característico
de cristas estreitas e pontiagudas e vales largos e suaves, observado em
ondas oceânicas reais. Para um ponto de referência $(x_0, z_0)$ no plano
da água em repouso, a posição deslocada por uma única componente de onda
é dada por:
\begin{align}
x &= x_0 + Q A \hat{D}_x \cos(\vec{k} \cdot \vec{x}_0 - \omega t) \label{eq:gerstner-x}\\
z &= z_0 + Q A \hat{D}_z \cos(\vec{k} \cdot \vec{x}_0 - \omega t) \label{eq:gerstner-z}\\
y &= A \sin(\vec{k} \cdot \vec{x}_0 - \omega t) \label{eq:gerstner-y}
\end{align}
em que $A$ é a amplitude da onda, $\vec{k}$ é o vetor de onda (magnitude
$2\pi/\lambda$, com $\lambda$ o comprimento de onda, e direção $\hat{D}$
de propagação), $\omega$ é a frequência angular, $t$ é o tempo, e
$Q \in [0,1]$ controla a nitidez das cristas (\emph{steepness}) ---
valores próximos de 1 produzem cristas mais pontiagudas, à custa de
possíveis autointerseções da malha caso o limite seja excedido. A
superfície final é obtida pela soma de $m = 6$ componentes de onda,
variando de ondulação longa a ondulação curta, com amplitudes
$A_i \in \{1{,}10;\, 0{,}90;\, 0{,}50;\, 0{,}28;\, 0{,}14;\, 0{,}07\}$ e
comprimentos de onda $\lambda_i \in \{70;\, 48;\, 26;\, 14;\, 7;\, 4\}$
unidades de mundo, com nitidez $Q_i$ decrescente de 0,60 a 0,28 à medida
que a frequência da componente aumenta --- restrição necessária para
evitar a autointerseção da malha nas componentes de alta frequência, cujo
comprimento de onda menor implica menor tolerância geométrica a cristas
pronunciadas. Essa superposição de componentes evita o padrão artificial
e repetitivo de uma única componente, aproximando visualmente o aspecto
irregular de uma superfície oceânica real, ao custo de $m$ avaliações da
Equação~\eqref{eq:gerstner-y} por vértice, a cada quadro renderizado.

A normal de cada ponto da superfície, necessária para a iluminação, é
obtida analiticamente a partir das derivadas parciais da soma das $m$
componentes de onda em relação a $x_0$ e $z_0$ --- e não por diferenças
finitas ---, o que evita o custo adicional de amostragem de vértices
vizinhos. O modelo de iluminação da superfície combina um termo difuso
lambertiano, um termo especular de Blinn-Phong, um termo de reflectância
dependente do ângulo de visão via aproximação de Fresnel-Schlick, e
transparência com fator $\alpha = 0{,}80$, o que permite a percepção
parcial do relevo submerso através da superfície da água.

A altura de base $y_0$ da superfície --- o plano de repouso ao redor do
qual as componentes de Gerstner oscilam --- é tratada como parâmetro
ajustável do sistema, controlável tanto de forma discreta, por meio de
quatro cenários pré-definidos, correspondentes a 0, +0,5, +1 e +2 metros
sobre o nível de referência atual, quanto de forma contínua, a uma taxa
de 0,8 m/s enquanto o comando correspondente é mantido pressionado. A
submersão de um ponto do terreno é determinada pela comparação direta
entre a elevação do terreno naquela posição (Equação~\eqref{eq:escala}) e
$y_0$: pontos cuja elevação é inferior a $y_0$ ficam abaixo da superfície
da água e, portanto, visualmente submersos, sem qualquer alteração na
geometria do terreno em si.

\subsection{Ambiente de Desenvolvimento}

A aplicação foi desenvolvida em linguagem Python, versão 3.11. A
biblioteca GLFW é responsável pela criação da janela e pela captura de
eventos brutos de teclado e mouse junto ao sistema operacional; a
biblioteca ModernGL fornece uma interface mais direta à API OpenGL
(\emph{core profile}, versão 3.3, correspondente à versão 330 da
linguagem de \emph{shading} GLSL) para a compilação de \emph{shaders}, a
alocação de \emph{Vertex Buffer Objects} (VBOs) e \emph{Vertex Array
Objects} (VAOs) na GPU, e a emissão de comandos de desenho (\emph{draw
calls}); a biblioteca PyGLM fornece os tipos e operações de álgebra
linear necessários à avaliação das Equações~\eqref{eq:front}
a~\eqref{eq:gerstner-y}; a biblioteca Pillow é utilizada tanto para a
decodificação dos \emph{tiles} de elevação (Equação~\eqref{eq:terrarium})
quanto para a renderização de texto da interface de usuário, uma vez que
a API OpenGL não oferece suporte nativo a texto; e a biblioteca
matplotlib é utilizada exclusivamente pelo módulo de avaliação
experimental (Seção IV) para a geração dos gráficos de desempenho.
Notavelmente, a leitura do modelo digital de elevação prescinde de
bibliotecas geoespaciais especializadas, como GDAL ou \texttt{rasterio}:
por tratar-se de \emph{tiles} no formato \emph{terrarium} --- imagens PNG
convencionais com a elevação codificada em canais de cor, conforme a
Equação~\eqref{eq:terrarium} ---, a decodificação é realizada diretamente
com Pillow e NumPy.

Os cálculos de seleção de nível da quadtree (Subseção III-B) e de
deslocamento de vértice das ondas de Gerstner (Subseção III-D) são
executados no estágio de \emph{vertex shader}, e não na CPU em Python, de
modo que o paralelismo massivo da GPU absorva o custo por vértice dessas
operações. O código-fonte está organizado em oito módulos principais:
\texttt{dem.py} (aquisição e pós-processamento do DEM real),
\texttt{heightmap.py} (geração procedural e grade unitária reusável pelo
CDLOD), \texttt{lod.py} (quadtree e seleção de blocos por quadro),
\texttt{water.py} (geometria do plano de água), \texttt{camera.py}
(câmera livre), \texttt{scene.py} (montagem da cena na GPU e desenho de
um quadro, reutilizada tanto pela aplicação interativa quanto pelo
\emph{benchmark}), \texttt{hud.py} (interface de usuário) e
\texttt{main.py} (laço principal), além do módulo \texttt{bench.py},
descrito na Seção IV. A interface de usuário inclui uma tela introdutória
com a descrição do objetivo do sistema, uma legenda de comandos em dois
níveis de detalhamento, um painel de status com o nível do mar corrente,
o modo de renderização ativo e o tempo de quadro medido, além de uma
bússola e um minimapa de orientação --- elementos voltados à
acessibilidade da ferramenta para um público não especializado, conforme
a motivação apresentada na Seção I.

A Tabela \ref{tab:parametros} sintetiza os parâmetros de configuração da
implementação de referência descrita nesta seção.

\begin{table}[t]
\caption{Parâmetros de configuração da implementação de referência}
\label{tab:parametros}
\centering
\begin{tabular}{@{}p{3.1cm}p{0.9cm}p{2.4cm}@{}}
\toprule
\textbf{Parâmetro} & \textbf{Símbolo} & \textbf{Valor} \\
\midrule
Resolução do \emph{heightmap} & $n$ & $1024 \times 1024$ \\
Exagero vertical & $E$ & $6\times$ \\
Escala do mundo & $\sigma$ & 167 m/u.m. \\
Profundidade máx. da quadtree & $L$ & 5 níveis \\
Resolução de vértices por nó & $N$ & $33 \times 33$ \\
Fator do critério de subdivisão & --- & $2{,}2 \times \tau$ \\
Componentes de onda de Gerstner & $m$ & 6 \\
Nitidez das ondas (faixa) & $Q$ & 0,60 a 0,28 \\
Transparência da água & $\alpha$ & 0,80 \\
Campo de visão da câmera & FOV & $60^{\circ}$ \\
Planos de recorte (próximo/distante) & --- & 0,1 / 2000 u.m. \\
Velocidade normal / turbo & --- & 40 / 120 u.m./s \\
\bottomrule
\end{tabular}
\end{table}

\section{Resultados e Avaliação de Desempenho}

A avaliação de desempenho foi conduzida por um programa de teste
(\emph{benchmark}) dedicado (\texttt{bench.py}), executado sem abrir
janela (modo \emph{headless}), reutilizando a mesma lógica de
montagem de cena (\texttt{scene.py}) empregada pela aplicação interativa,
de modo a garantir que a medição reflita o custo real de renderização, e
não uma via de código simplificada. O protocolo experimental sincronizou
a GPU ao final de cada quadro, evitando que o \emph{buffering} assíncrono
mascarasse o tempo real de processamento; descartou 30 quadros iniciais
de aquecimento (\emph{warm-up}); e mediu os 150 quadros subsequentes. O
experimento variou três fatores: (i) a técnica de renderização --- malha
única na resolução total (sem LOD), CDLOD completo, e CDLOD sem a etapa
de \emph{geomorphing} (Subseção III-B); (ii) a resolução do
\emph{heightmap}, em cinco níveis (256, 512, 1024, 1536 e 2048); e (iii)
a posição da câmera, em quatro configurações distintas dentro da cena.
Adicionalmente, foi registrado o consumo de memória do processo (RAM) e
o uso de memória de vídeo e de GPU (VRAM), por meio das bibliotecas
\texttt{psutil} e \texttt{nvidia-smi}, respectivamente. Os testes foram
executados em uma GPU NVIDIA RTX 3060.

A Figura~\ref{fig:fps} apresenta a taxa de quadros média (das quatro
posições de câmera) em função da resolução do \emph{heightmap}, para as
três condições avaliadas.

\begin{figure}[t]
\centering
\includegraphics[width=\columnwidth]{fps_vs_resolucao.png}
\caption{Taxa de quadros média em função da resolução do \emph{heightmap}:
o CDLOD (com e sem \emph{geomorphing}) mantém custo de renderização
praticamente constante, enquanto a malha sem LOD degrada linearmente com
o aumento da resolução.}
\label{fig:fps}
\end{figure}

Os resultados confirmam a hipótese subjacente ao CDLOD: como o algoritmo
restringe a contagem de blocos ativos a um valor aproximadamente
constante (200 a 230 blocos, Subseção III-B), independentemente da
resolução total do \emph{heightmap}, o tempo de quadro da condição CDLOD
permanece próximo de 1 ms em todas as resoluções testadas,
correspondendo a uma carga de renderização estável de aproximadamente
0,4 milhão de triângulos por quadro. Em contraste, a condição sem LOD ---
que renderiza a malha completa a cada quadro --- apresenta tempo de
quadro crescente linearmente com o número de triângulos da malha,
chegando a aproximadamente 8,4 milhões de triângulos na resolução máxima
testada (2048), frente a apenas 2,1 milhões de triângulos na resolução
intermediária de 1024.

As curvas de desempenho das duas condições cruzam-se em aproximadamente
cinco milhões de triângulos de carga geométrica --- o ponto de virada
(\emph{tipping point}) a partir do qual o custo da técnica CDLOD passa a
ser inferior ao da renderização da malha completa nessa GPU. Abaixo desse
ponto, a malha sem LOD é, de fato, ligeiramente mais rápida, uma vez que
o CDLOD introduz uma sobrecarga fixa de gerenciamento da quadtree que não
se paga para cargas geométricas pequenas. Esse resultado constitui a
validação direta da hipótese enunciada na Seção I: o CDLOD não é
vantajoso de forma incondicional, mas passa a ser indispensável a partir
de uma carga geométrica identificável experimentalmente --- e, em uma GPU
de menor capacidade de processamento, como as de dispositivos móveis
(público-alvo final da aplicação), espera-se que esse ponto de virada
ocorra em uma resolução de terreno ainda menor, tornando o CDLOD
relativamente mais vantajoso nesse contexto. O custo computacional
específico da etapa de \emph{geomorphing} (comparação entre as condições
``CDLOD'' e ``CDLOD sem \emph{morphing}'' na Figura~\ref{fig:fps})
revelou-se dentro da margem de ruído da medição, isto é, sua eliminação
de \emph{cracks} e \emph{popping} (Subseção III-B) tem, na prática, custo
de desempenho desprezível.

Quanto ao consumo de recursos, observou-se crescimento da memória do
processo (RAM) de aproximadamente 110 MB, na resolução 256, a 255 MB, na
resolução 2048, e da memória de vídeo (VRAM) de aproximadamente 740 MB a
910 MB no mesmo intervalo. Esse crescimento acompanha a resolução do
\emph{dado de entrada} (o tamanho da textura de altura e dos
\emph{buffers} correspondentes), e não o modo de renderização
selecionado: nas condições avaliadas, a malha completa e as duas
variantes do CDLOD compartilham, dentro de uma mesma execução, os
mesmos recursos residentes na GPU, de modo que este experimento não
isola o ganho de memória atribuível especificamente ao LOD, tarefa
identificada como trabalho futuro (Seção V).

Verificou-se, por fim, que a forma do terreno gerado --- os canais do
sistema estuarino de Santos e São Vicente, as ilhas próximas e a
orientação da escarpa da Serra do Mar --- corresponde à geografia real da
região quando comparada a um mapa de referência, confirmando a
integridade do processo de aquisição e decodificação descrito na
Subseção III-A.

\section{Conclusão e Trabalhos Futuros}

Este artigo apresentou um sistema interativo que exibe em 3D o relevo
real da Baixada Santista e uma superfície de água animada, cuja altura
pode ser elevada em 0, $+0{,}5$, $+1$ e $+2$ metros, permitindo
observar quais áreas ficariam submersas em cada cenário de subida do
mar.

O principal resultado vem do teste de desempenho. Com o CDLOD, o tempo
para desenhar um quadro da cena completa (terreno e água) ficou próximo
de 1 ms em todas as resoluções de terreno testadas. Sem LOD, esse tempo
cresce junto com o tamanho da malha. As duas curvas se cruzam em
aproximadamente cinco milhões de triângulos: abaixo desse valor, desenhar
a malha completa é ligeiramente mais rápido; acima dele, o CDLOD passa a
ser indispensável para manter a interatividade. O custo do
\emph{geomorphing} mostrou-se desprezível. A eficácia do CDLOD já é
conhecida na literatura \cite{strugar2009}; o que este trabalho
acrescenta é a medição desse ponto de virada num cenário aplicado de
inundação costeira.

Algumas limitações devem ser consideradas. Primeiro, nossa implementação
do CDLOD é simplificada: não descarta os blocos que estão fora do campo
de visão da câmera (\emph{frustum culling}) e as faixas de transição do
\emph{geomorphing} não seguem exatamente as fórmulas do trabalho original
\cite{strugar2009}. Segundo, o teste foi feito apenas em uma GPU de
computador de mesa (RTX 3060); a expectativa de que o ponto de virada
ocorra mais cedo em celulares, nosso público-alvo final, é uma hipótese
razoável, mas não foi medida. Terceiro, o relevo foi exagerado
verticalmente em $6\times$ (Equação~\eqref{eq:escala}) para que a Serra
do Mar fosse perceptível na escala da área; trata-se de uma distorção
deliberada, comum em geovisualização, mas que precisa ser declarada.
Quarto, a água não interage com o terreno (não há espuma na costa nem
refração da luz), e a medição de memória não separa o efeito do LOD,
pois os três modos testados compartilham os mesmos recursos na GPU.

Como trabalhos futuros, propomos quatro passos, cada um respondendo a
uma limitação acima. (i) Repetir o teste de desempenho em um celular ou
em uma GPU mais fraca, para confirmar o ponto de virada nesses
dispositivos e medir a memória de cada modo separadamente. (ii)
Implementar o \emph{frustum culling}, de modo a não desenhar os blocos
atrás da câmera e reduzir ainda mais o custo por quadro. (iii) Acrescentar
espuma e arrebentação na linha de costa, para tornar a água mais
realista. (iv) Apresentar a ferramenta a especialistas em gestão
costeira, a fim de avaliar se ela de fato ajuda a comunicar o risco de
inundação.

\begin{thebibliography}{00}

\bibitem{tanner1998} C.~C. Tanner, C.~J. Migdal, e M.~T. Jones, ``The
clipmap: a virtual mipmap,'' in \emph{Proceedings of the 25th Annual
Conference on Computer Graphics and Interactive Techniques (SIGGRAPH)},
1998, pp. 151--158.

\bibitem{clasenhege2006} M.~Clasen e H.-C. Hege, ``Terrain rendering
using spherical clipmaps,'' in \emph{Proceedings of EuroVis 2006}, 2006,
pp. 91--98.

\bibitem{strugar2009} F.~Strugar, ``Continuous distance-dependent level
of detail for rendering heightmaps,'' \emph{Journal of Graphics, GPU,
and Game Tools}, vol. 14, no. 4, pp. 57--74, 2009.

\bibitem{fournierreeves1986} A.~Fournier e W.~T. Reeves, ``A simple model
of ocean waves,'' \emph{ACM SIGGRAPH Computer Graphics}, vol. 20, no. 4,
pp. 75--84, 1986.

\bibitem{welch} R.~Welch, ``Real-time ocean simulation,'' trabalho de
conclusão de curso (B.Sc. em Ciência da Computação), 2019. Disponível em:
\url{https://ryanwelch.co.uk/project/ocean-simulation/}.

\bibitem{brown2006} I.~Brown, S.~Jude, S.~Koukoulas, R.~Nicholls,
M.~Dickson, e M.~Walkden, ``Dynamic simulation and visualisation of
coastal erosion,'' \emph{Computers, Environment and Urban Systems},
vol. 30, no. 6, pp. 840--860, 2006.

\bibitem{dipaola2023} F.~DiPaola, A.~Bhardwaj, e L.~Sam, ``Generating an
interactive online map of future sea level rise along the North Shore of
Vancouver: methods and insights on enabling geovisualisation for coastal
communities,'' \emph{arXiv preprint arXiv:2304.07469}, 2023.

\end{thebibliography}

\end{document}
